% ATTEMPT_STATUS: unresolved
% TOP_PROBLEM: {"problemId": "problem.the-mumford-tate-conjecture-for-abelian-varieties-over-number-fields", "canonicalTitle": "The Mumford-Tate Conjecture for Abelian Varieties over Number Fields", "releaseRank": 92, "primaryDomain": "number_theory_arithmetic_geometry", "primaryDomainLabel": "Number theory & arithmetic geometry", "displayStatus": "open", "releaseStatus": "open"}
% !TeX program = lualatex
\documentclass[11pt]{article}
\usepackage[margin=25mm]{geometry}
\usepackage{fontspec}
\setmainfont{Latin Modern Roman}
\usepackage{amsmath,amssymb,mathtools}
\usepackage{unicode-math}
\setmathfont{Latin Modern Math}
\usepackage{xurl,hyperref}
\hypersetup{colorlinks=true,urlcolor=blue}
\setcounter{secnumdepth}{0}
\setlength{\parindent}{0pt}
\setlength{\parskip}{5pt}
\setlength{\emergencystretch}{3em}
\begin{document}
\section*{\texorpdfstring{The Mumford-Tate Conjecture for Abelian Varieties over Number Fields}{The Mumford-Tate Conjecture for Abelian Varieties over Number Fields}}\label{92}
\subsection{Definitions and mathematical statement}
For an abelian variety $A/K$, let $V_\ell(A)=T_\ell(A)\otimes{\mathbb{Q}}_\ell$ and let $G_\ell$ be the Zariski closure of the Galois image in $\operatorname{GL}(V_\ell(A))$. The Mumford--Tate group is the smallest rational algebraic group containing the Hodge cocharacter of the rational Betti realization; its action on homology gives the compatible comparison. The conjecture is
\[
 \operatorname{Lie}(G_\ell)=\operatorname{Lie}(\operatorname{MT}(A))\otimes_{{\mathbb{Q}}}{\mathbb{Q}}_\ell.
\]
Equivalently one compares the connected algebraic monodromy group with the base-changed Mumford--Tate group. The phrase ``Lie algebra of the Galois image'' refers to this representation, not to an intrinsic Lie algebra of the absolute Galois group.
\par\smallskip\textit{Formulation and concept references:} \href{https://doi.org/10.1201/9781439864227}{[S1]}.

\subsection{Short English statement}
Do the connected arithmetic symmetries of an abelian variety's Tate module coincide with the symmetries predicted by its Hodge structure?
\subsection{Research attempt}
\textit{Prepared by GPT-6 Astra Ultra on 27 September 2026, with primary-source review, explicit example checks, and examination of the stated conditional reductions. This is an unresolved research attempt.}

\paragraph{Scope and source check (27 September 2026).}
Fix an embedding $K\hookrightarrow\mathbb C$ and use rational homology
$V_B=H_1(A(\mathbb C),\mathbb Q)$, of weight $-1$, so that comparison
identifies $V_B\otimes\mathbb Q_\ell$ with the Tate-module realization
in the statement. All groups below are algebraic groups in these
representations. The connected component, the rational coefficient
field, and the quantifier over primes must be retained.

The strategy is to test whether the known endomorphism comparison
already forces equality, then replace that insufficient test by a
precise tensor criterion. Pink [E1, Section 5] records the inclusion
of connected monodromy in the Mumford--Tate group and the consequences
of Faltings' endomorphism theorem. These are established inputs, not
results proved here. Source [S2, Theorem 4.16] also uses that inclusion;
its Frobenius-field density theorem does not establish the reverse
inclusion. The calculation below recovers a known elliptic-curve case
and constructs an explicit obstruction to an overly short general proof.

\paragraph{Finite extension does not change the target.}
Let $K'/K$ be finite. The corresponding Galois images $\Gamma'\subset
\Gamma$ have finite index, so, on taking Zariski closures,
\[
 G_\ell=\bigcup_{j=1}^t g_jG'_\ell.
\]
Consequently $G'_\ell$ has finite index as an algebraic subgroup and
$(G'_\ell)^\circ=G_\ell^\circ$: the quotient into the finite set of
cosets is constant on a connected component. This can also be checked
after passing to the finite-index normal core. The complex abelian
variety and its rational Hodge structure have not changed. Thus a
finite extension preserves exactly the desired equality of connected
groups.

For a fixed $\ell$ one may therefore arrange that the monodromy group
is connected and that every geometric endomorphism is defined over
the ground field. The latter needs only finitely many extensions:
the endomorphism ring is finitely generated over $\mathbb Z$, and each
chosen generator is defined over a finite extension. Connectedness is
obtained by taking the preimage of the identity component in the
continuous Galois representation. This observation does not require
one common extension for all primes.

\paragraph{What the established comparison actually supplies.}
Put $G=\operatorname{MT}(A)_{\mathbb Q_\ell}$ and $H=G_\ell^\circ$.
The absolute-Hodge-cycle inclusion and Faltings' theorem give
\begin{equation}\label{mt:known}
 H\subseteq G,\qquad
 \operatorname{End}_H(V_\ell)
 =\operatorname{End}(A_{\overline K})\otimes\mathbb Q_\ell
 =\operatorname{End}_G(V_\ell).
\end{equation}
For the second equality one uses the usual full faithfulness of rational
homology for complex abelian varieties, and the fact that their
endomorphisms descend to $\overline K$. Homology makes the action
covariant; the opposite algebras appearing in a cohomology convention
are not needed here. Both groups are reductive: on the arithmetic
side this follows from semisimplicity, and on the Hodge side from
polarizability. None of these statements identifies the groups merely
from their endomorphism rings.

A polarization places both in a group of symplectic \emph{similitudes}.
It does not place the full groups in a symplectic group: the Tate line
and the homotheties carry the multiplier. Similarly, absolute Hodge
cycles are not being asserted to be algebraic cycles. The established
inclusion in \eqref{mt:known} uses their absolute-Hodge property.

\paragraph{An explicit countertest for the commutant argument.}
Over $\mathbb Q$ take $W=\operatorname{Sym}^3(\mathbb Q^2)$ with basis
$e_i=x^{3-i}y^i$, $0\leq i\leq3$. Define an alternating form $J$ by
\[
 J(e_i,e_j)=0\quad(i+j\ne3),\qquad
 J(e_i,e_{3-i})=\frac{(-1)^i}{\binom3i}.
\]
The four antidiagonal entries are $1,-1/3,1/3,-1$, so the form is
nondegenerate. The infinitesimal operators
\[
 E e_i=i e_{i-1},\qquad F e_i=(3-i)e_{i+1},\qquad
 D e_i=(3-2i)e_i
\]
preserve $J$ infinitesimally. For example, the only potentially nonzero
terms in $J(Ee_i,e_j)+J(e_i,Ee_j)$ have $i+j=4$ and cancel by
$i/\binom3{i-1}=j/\binom3i$; the signs are opposite. The analogous
calculation works for $F$, and the paired $D$-weights sum to zero.
Thus the connected group
\[
 H_0=\mathbb G_m\cdot\operatorname{Sym}^3(\operatorname{SL}_2)
 \subset \operatorname{GSp}(W,J)=G_0
\]
is reductive and proper: its dimension is $4$, whereas $\dim G_0=11$.
It contains all homotheties and its multiplier has image $\mathbb G_m$
as an algebraic group.

Nevertheless both commutants are the scalars. Indeed an endomorphism
commuting with $D$ is diagonal because its four weights are distinct;
commuting with $E$ forces its successive diagonal entries to agree.
Hence $\operatorname{End}_{H_0}(W)=\mathbb Q$, and the same follows
for $G_0$ by inclusion. The reasoning is unchanged after field
extension. This explicitly disproves the abstract inference that
equal commutants, reductivity, a polarization, and scalar matrices
force equality.

This is not a counterexample to the conjecture. No abelian variety has
been produced with these two groups. In particular the restricted
Hodge-cocharacter conditions on abelian-variety monodromy are extra
information, and this example has not been shown to satisfy them.
Its role is to locate exactly what the endomorphism-only argument omits.

\paragraph{A sufficient replacement using all mixed tensors.}
For a characteristic-zero field $k$ and $V$ over $k$, write
\[
 T^{a,b}(V)=V^{\otimes a}\otimes(V^*)^{\otimes b}.
\]
If $H\subseteq G\subseteq\operatorname{GL}(V)$ are reductive algebraic
groups, then
\begin{equation}\label{mt:tensors}
 H=G\quad\Longleftrightarrow\quad
 T^{a,b}(V)^H=T^{a,b}(V)^G\quad\hbox{for all }a,b\geq0.
\end{equation}
Here equality means equality of the actual subspaces under the same
representation, not just an unexplained comparison of dimensions.
In the present nested situation equality of dimensions would imply
equality of subspaces, since $G$-invariants are already $H$-invariants.

For completeness, the reverse implication follows from the standard
Chevalley line-stabilizer theorem [E2]. Realize $H$ as the stabilizer
of a line $L$ in a finite-dimensional rational representation $U$ of
$\operatorname{GL}(V)$. As $H$ is reductive, choose an $H$-stable
complement to $L$. The projector $P:U\to L$ then commutes with $H$.
Every such $U$ is a direct summand of a finite direct sum of mixed
tensor powers: polynomial representations together with determinant
inverses suffice, and characteristic-zero complete reducibility
turns their subquotients into summands. Consequently the hypothesized
tensor equalities apply also to $\operatorname{End}(U)$. They force
$P$ to commute with $G$, so $G$ preserves its image $L$. By the choice
of $L$, this gives $G\subseteq H$, proving equality.

Applied to \eqref{mt:known}, this gives a precise conditional route.
One must prove that every $H$-fixed mixed tensor belongs to the
$\mathbb Q_\ell$-span of the rational Hodge tensors fixed by
$\operatorname{MT}(A)$. An arbitrary $\ell$-adic invariant need not
itself have rational coordinates, so that span formulation is
essential. The reverse containment is already known. Endomorphism
comparison supplies just $T^{1,1}$; the line projector in the proof
can occur much farther up the tensor algebra. For each hypothetical
proper subgroup there is some finite distinguishing tensor, but this
existence proof supplies neither a uniform small degree nor an
arithmetic method for ruling out all such tensors. Replacing this
missing step by the phrase ``Tannakian formalism'' would only rename it.

\paragraph{A complete restricted case: a non-CM elliptic curve.}
Suppose $A=E$ is an elliptic curve with
$\operatorname{End}(E_{\overline K})=\mathbb Z$. The inputs above imply
that both $H$ and $G$ act absolutely irreducibly on their two-dimensional
standard representations. Indeed their scalar commutants remain
one-dimensional after algebraic closure; a reducible semisimple
representation would have at least two independent commuting
projections.

Over an algebraically closed field of characteristic zero, a connected
reductive irreducible subgroup of $\operatorname{GL}_2$ has derived
group $\operatorname{SL}_2$. If its derived Lie algebra were zero, the
group would be a torus and its representation would split into lines.
Otherwise that Lie algebra is a nonzero semisimple subalgebra of
$\mathfrak{sl}_2$, and hence is all $\mathfrak{sl}_2$: a proper
subalgebra has dimension at most two and is solvable. Connectedness
then gives the asserted derived group.

Finally $\det H=\mathbb G_m$. On a Tate module the determinant is the
$\ell$-adic cyclotomic character; its image over a number field, even
after finite extension, is infinite and therefore Zariski dense in
$\mathbb G_m$. Similarly $\det G=\mathbb G_m$, since the weight
homotheties lie in the Mumford--Tate group. A subgroup containing
$\operatorname{SL}_2$ and having this determinant image is
$\operatorname{GL}_2$. Equality over the algebraic closure descends
to the original fields, so $H=G=\operatorname{GL}_{2,\mathbb Q_\ell}$.
This recovers a classical known subcase, using deep established
inputs; it is not a new proof of the general conjecture.

\paragraph{Why Frobenius spectra alone do not finish the route.}
The closure of the cyclic subgroup generated by a semisimple
Frobenius matrix has torus identity component. Even when that is a
maximal torus, its eigenvalues describe weights, not all root
subgroups. For example a torus in $H_0$ above has eigenvalues
$\lambda t^3,\lambda t,\lambda t^{-1},\lambda t^{-3}$, while $H_0$
is still a proper similitude subgroup. This is not a claim that its
formal character equals that of every possible ambient group.
It shows why finding an interesting Frobenius torus is insufficient
without further group structure and arithmetic information.
Nor does a calculation at one prime automatically establish the
conjecture at every prime without a separately justified independence
theorem in the applicable scope.

\paragraph{Verification, first gap, and outcome.}
Exact rational matrix checks verify the displayed alternating form,
the three infinitesimal invariance identities, and the one-dimensional
commutant in the four-dimensional countertest. The group dimensions,
finite-extension reduction, tensor criterion, and elliptic case were
also checked at the level of their stated hypotheses. These are tests
of the deductions, not computations of general Galois images.
The first missing general input is the reverse inclusion on higher
tensor invariants, or an equally strong classification argument
excluding every proper arithmetic subgroup allowed by the Hodge
data. Neither has been established here.
\textbf{Outcome: UNRESOLVED.} The attempt provides a rigorous restricted
case and a concrete failure test for a plausible shortcut, not a
general proof or a claim of novelty.

\subsection{Sources}
\begin{itemize}
\item[S1]J.-P. Serre, Abelian l-Adic Representations and Elliptic Curves, Addison-Wesley, 1968.
\url{https://doi.org/10.1201/9781439864227}.
\textit{Formulation source.}
\item[S2]On the Frobenius fields of abelian varieties over number fields.
\url{https://arxiv.org/html/2402.07935v2}.
\textit{Further reference; not a proof endorsement.}
\item[E1]Richard Pink, \emph{$\ell$-adic algebraic monodromy groups, cocharacters, and the Mumford--Tate conjecture}, Journal f\"ur die reine und angewandte Mathematik 495 (1998), 187--237; author manuscript dated 14 May 1997.
\url{https://people.math.ethz.ch/~pinkri/ftp/AMT-v3.pdf}.
\textit{Section 5 records the established inclusion, endomorphism comparison, and representation restrictions. Consulted 27 September 2026. The elementary countertest and restricted deduction above are supplied explicitly.}
\item[E2]J. S. Milne, \emph{Algebraic Groups}, corrected reprint, Cambridge University Press, 2022, Chapter 4, especially Theorem 4.27.
\url{https://www.jmilne.org/math/Books/iAG2022.pdf}.
\textit{Author text for rational representations and Chevalley's line-stabilizer theorem; consulted 27 September 2026.}
\end{itemize}
\end{document}
